\documentclass[11pt]{article}

\usepackage[utf8]{inputenc}
\usepackage[T1]{fontenc}
\usepackage{lmodern}
\usepackage{geometry}
\usepackage{amsmath,amssymb,amsfonts,amsthm,mathtools}
\usepackage{bm}
\usepackage{enumitem}
\usepackage{microtype}
\usepackage{hyperref}
\usepackage{titlesec}
\usepackage{booktabs}
\usepackage{array}
\usepackage{setspace}
\usepackage{mathrsfs}
\usepackage{physics}

\geometry{margin=1in}
\hypersetup{colorlinks=true, linkcolor=black, urlcolor=blue, citecolor=black}
\setlist[enumerate]{topsep=0.4em,itemsep=0.35em}
\setlist[itemize]{topsep=0.3em,itemsep=0.25em}
\titleformat{\section}{\large\bfseries}{\thesection.}{0.5em}{}
\titleformat{\subsection}{\normalsize\bfseries}{\thesubsection.}{0.5em}{}
\setstretch{1.06}

\newcommand{\Aether}{\AE{}ther}
\newcommand{\Aflow}{\AE{}ther-flow}
\newcommand{\TheoryName}{\Aflow{} Interpretation of Relativity}
\newcommand{\TheoryProgram}{\Aether{} / \Aflow{} framework}

\newtheorem{definition}{Definition}
\newtheorem{proposition}{Proposition}
\newtheorem{remark}{Remark}
\newtheorem{corollary}{Corollary}

\title{\textbf{The \TheoryName{}}\\[0.4em]
\large Exact-Closure Sequence Overview}
\author{}
\date{}

\begin{document}

\maketitle

\begin{abstract}
This note presents the active exact-closure manuscript sequence for \TheoryName{} as the flagship public-facing deliverable of the repository. Its purpose is not to replace the individual root manuscripts, but to fix their common role, reading order, and shared claim boundary while packaging them as one coherent theory statement. In the active repository state, \TheoryName{} is the exact relativistic theory built on the \Aether{} / \Aflow{} ontology. Its effective gravitational dynamics are adopted to be exactly those of general relativity with universal matter coupling. Accordingly, the exact-closure sequence already provides a disciplined theory statement at the effective level even though the deeper substrate derivation remains open.

The overview makes four points explicit. First, this note is the canonical front door to the flagship exact-theory package. Second, the exact-closure note, foundations, dynamics, consistency, relativistic-recovery, and flow-geometry papers should be read as one fixed coordinated line after that front door. Third, the predictive content of the active theory is exactly that of GR, while the \Aether{} / \Aflow{} vocabulary remains ontological and interpretive at the effective level. Fourth, the derivational bridge manuscripts belong to an ongoing foundational recovery program answerable to this exact-closure benchmark rather than to the completed theory statement itself.
\end{abstract}

\section{Introduction}

The active repository has reached a stage at which the exact-closure line should be presented as the flagship scientific deliverable in its own right. The deeper framework still permits ongoing derivational work, but the project is no longer best described as ontology awaiting its first disciplined theory statement. That theory statement now exists in the exact-closure sequence.

\paragraph{Framework and claim boundary.}
\TheoryProgram{} denotes the broader \Aether{} / \Aflow{} framework built on the claims that the \Aether{} is the underlying four-dimensional substrate of reality, the \Aflow{} is its intrinsic ordered motion, observed three-dimensional space is the local experiential slice of that deeper substrate, S-time is the experienced order of change arising from matter, light, and the \Aflow{}, and observed expansion is the three-dimensional appearance of deeper four-dimensional ordered motion. Within that framework, \TheoryName{} denotes the exact-closure theory in which the effective gravitational dynamics are adopted to be exactly Einsteinian with universal matter coupling. In the active sequence, \emph{adoption} means use of that established relativistic dynamics without claiming substrate derivation, while \emph{derivation} is reserved for a first-principles recovery from explicit substrate variables.

This distinction is the organizing principle of the overview. The active exact-closure line is already a positive deliverable. The still-open substrate program is a later foundational burden, not the condition for calling the effective theory statement real. The immediate packaging task is therefore consolidation of the flagship exact-theory sequence rather than another derivational bridge installment.

\section{Purpose of the Overview}

\begin{definition}[Exact-closure sequence overview]
An exact-closure sequence overview is a manuscript that presents the active root papers of \TheoryName{} as one coordinated theory statement, identifies the role of each paper within that statement, and fixes the boundary between completed exact closure and unfinished substrate derivation.
\end{definition}

\begin{proposition}[Current architectural status]
At the present stage of the repository, the canonical public-facing entry to \TheoryProgram{} is the exact-closure manuscript sequence rather than the derivational bridge program.
\end{proposition}

This proposition is not merely editorial. It reflects the scientific order now achieved in the repository. The exact operational theory, its consistency status, and its relation to relativistic phenomenology are already stated in disciplined manuscript form. What remains open is not whether the project has any exact theory statement, but whether that exact theory can later be recovered from a more primitive substrate dynamics.

\begin{corollary}[Flagship packaging rule]
The fixed flagship package of the repository is this overview together with the six core exact-closure manuscripts, and any further derivational continuation must be presented only as downstream foundational work.
\end{corollary}

\section{Exact Theory Statement}

\begin{definition}[\TheoryName{}]
\TheoryName{} is the exact-closure, GR-consistent theory of \TheoryProgram{} in which the \Aether{} / \Aflow{} ontology is retained while the effective gravitational dynamics are exactly those of general relativity with universal matter coupling.
\end{definition}

Its effective action is
\begin{equation}
S_{\mathrm{eff}}
=
\frac{c^3}{16\pi G}\int d^4x\,\sqrt{-g}\,R
+
S_{\mathrm{matter}}[g,\psi],
\label{eq:SA}
\end{equation}
with field equations
\begin{equation}
G_{\mu\nu}
=
\frac{8\pi G}{c^4}T_{\mu\nu},
\label{eq:einstein}
\end{equation}
the standard null condition
\begin{equation}
0=g_{\mu\nu}\,dx^\mu dx^\nu,
\label{eq:null}
\end{equation}
and the standard proper-time rule
\begin{equation}
d\tau^2=-\frac{1}{c^2}g_{\mu\nu}\,dx^\mu dx^\nu.
\label{eq:propertime}
\end{equation}

\begin{proposition}[Operational identity]
For fixed matter sector, initial data, and boundary conditions, the effective observable content of \TheoryName{} is identical to that of GR.
\end{proposition}

\begin{remark}
This is an adoption statement, not a first-principles derivation statement. The exact-closure sequence does not claim that the Einstein sector has already been recovered uniquely from deeper substrate microphysics.
\end{remark}

\section{Flagship Package and Reading Order}

This note is the canonical front door to the exact-closure sequence. After it, the active core manuscripts should be read in the following fixed order and should be presented in this same order in repository-facing summaries.
\begin{enumerate}
    \item \textbf{Exact-closure note.} The standalone exact-closure note states the shortest complete version of the active scientific position, fixes the benchmark role of \TheoryName{}, and marks the nonclaims that preserve discipline.
    \item \textbf{Foundations.} The foundations manuscript fixes the ontological vocabulary, clarifies the original intent of the framework, and states why exact closure rather than unfinished deformation is the correct active stance.
    \item \textbf{Dynamics.} The dynamics manuscript states the adopted effective action, field equations, matter-coupling rule, weak-field structure, and benchmark role of the exact-closure theory.
    \item \textbf{Consistency.} The consistency manuscript makes explicit the gauge structure, degree-of-freedom count, and effective health conditions inherited from GR in the adopted exact-closure theory.
    \item \textbf{Relativistic recovery.} The relativistic-recovery manuscript states the exact relation to GR and the local relation to SR, including weak-field recovery, null propagation, redshift, clock behavior, and local inertial structure.
    \item \textbf{Flow geometry.} The flow-geometry manuscript supplies the disciplined congruence-based dictionary by which the \Aflow{} ontology is read inside exact relativistic geometry without introducing a second low-energy law.
\end{enumerate}

\begin{corollary}[Stable sequence logic]
Read in that order, the exact-closure papers supply one continuous argument: ontology is stated, exact dynamics are fixed, internal consistency is controlled, relativistic recovery is made explicit, and the \Aflow{} language is given formal geometric discipline.
\end{corollary}

This fixed reading order is part of the claim boundary of the flagship package. It prevents the effective theory statement from being blurred together with optional derivational continuation.

\section{What the Sequence Establishes}

The exact-closure sequence now establishes the following at the active theory level.
\begin{enumerate}
    \item The ontological vocabulary of \Aether{}, \Aflow{}, observed three-dimensional space, and S-time is fixed in disciplined form.
    \item The effective gravitational law is fixed exactly as Einsteinian with universal matter coupling.
    \item The active theory carries no independent low-energy non-GR observable signature.
    \item The effective gravitational sector has the ordinary massless spin-2 structure, standard relativistic causal order, and no additional established low-energy scalar or vector graviton sector.
    \item The observer-accessible weak-field, redshift, time-dilation, null-propagation, and local SR content is inherited exactly from the same metric structure.
    \item The \Aflow{} interpretation is formalized through congruence geometry rather than loose metaphor or a naive vector-wind picture.
\end{enumerate}

\begin{proposition}[Working theory standard]
At the effective level, \TheoryName{} already satisfies the working standard required for the project's main goal: ontology is matched by closed mathematical structure, effective dynamics are explicit, internal consistency is controlled, and observable relativistic structure is recovered without ambiguity.
\end{proposition}

\section{What the Sequence Does Not Establish}

The same discipline requires an explicit statement of what the consolidated exact-closure line does not yet claim.
\begin{enumerate}
    \item It does not claim a completed first-principles derivation of Einsteinian gravity from explicit substrate variables.
    \item It does not claim a unique underlying substrate action from which the exact relativistic sector has already been proved to emerge.
    \item It does not claim a distinct verified low-energy observable sector beyond GR.
    \item It does not claim that the broader bridge program has already solved uniqueness, mode control, or infrared emergence in full generality.
    \item It does not claim that ontology by itself constitutes empirical evidence for new low-energy physics.
\end{enumerate}

\begin{remark}
These nonclaims are not weaknesses of the overview. They are part of what makes the exact-closure sequence scientifically stable. The project now succeeds as an exact theory statement because it distinguishes completed effective structure from unfinished microphysical ambition.
\end{remark}

\section{Relation to the Derivational Bridge Program}

The broader substrate manuscript line remains important, but its role must be read correctly. It is answerable to the exact-closure benchmark already fixed above. Its purpose is to ask whether the exact relativistic sector can be recovered from deeper substrate variables while preserving the same operational content and without generating a surviving deformation of that content.

This means the bridge program is no longer the repository's only route to theory status. It is a secondary foundational research program built downstream of an already stated exact theory. That ordering matters. If a future bridge candidate fails, the scientific status of the project contracts back to \TheoryName{} rather than back to ontology alone. If the bridge line is deliberately resumed, it should be introduced explicitly as optional continuation beyond the flagship package rather than as the next required milestone.

\begin{proposition}[Benchmark and reversion role]
Any admissible future extension of \TheoryProgram{} must remain answerable to \TheoryName{} as exact benchmark, and failure of such an extension contracts the framework back to \TheoryName{} rather than invalidating the exact-closure sequence.
\end{proposition}

\section{Editorial Use of This Note}

This overview is the canonical front-door manuscript for the flagship exact-closure line and should be used whenever a short but complete statement of the active project status is needed. It is especially appropriate when the repository must be introduced at sequence level rather than through one root paper in isolation.

The note does not replace the six core manuscripts. Instead, it fixes how they fit together:
\begin{itemize}
    \item the exact-closure note remains the shortest standalone anchor;
    \item the five companion papers remain the full modular statement of foundations, dynamics, consistency, relativistic recovery, and flow geometry;
    \item the substrate notes remain an open derivational program downstream of the exact-closure benchmark and should be presented only after the flagship package is made explicit.
\end{itemize}

\section{Conclusion}

The main project goal is now best represented by the exact-closure sequence taken as one coordinated deliverable. \TheoryName{} is no longer merely an ontological suggestion awaiting its first disciplined theory statement. It is already the exact relativistic theory of the active framework at the effective level. What remains open is the deeper question of derivation from substrate first principles, not the existence of an exact operational theory.

Accordingly, the correct next architectural stance is consolidation rather than further blurring. The exact-closure sequence should be treated as the canonical flagship theory statement of the repository, while the derivational bridge program should continue only as a clearly subordinate and clearly labeled foundational research line.

\end{document}
